\begin{figure}[htbp]
\centering
\begingroup
\definecolor{ladderaccent}{RGB}{22,91,119}
\begin{tikzpicture}[x=1cm,y=1cm,>=Latex,
  every node/.style={font=\small},
  guide/.style={line width=.9pt},
  oldrung/.style={densely dashed,line width=.65pt,black!65},
  newrung/.style={line width=1.1pt,ladderaccent}]
  \coordinate (am) at (0,2.05);
  \coordinate (ai) at (5.4,2.05);
  \coordinate (ap) at (10.8,2.05);
  \coordinate (bm) at (0,.2);
  \coordinate (bo) at (6.7,.2);
  \coordinate (bn) at (5.4,.2);
  \coordinate (bp) at (10.8,.2);
  \draw[guide] (am)--(ap);
  \draw[guide] (bm)--(bp);
  \draw[oldrung] (am)--node[left] {$\le h$} (bm);
  \draw[oldrung] (ap)--node[right] {$\le h$} (bp);
  \draw[oldrung] (ai)--node[pos=.50,right,fill=white,inner sep=1.5pt]
    {old: $\le h$} (bo);
  \draw[newrung] (ai)--node[pos=.55,left,fill=white,inner sep=1.5pt]
    {new: $\le 2\Delta$} (bn);
  \foreach \p in {am,ai,ap,bm,bo,bp}
    \fill (\p) circle (1.7pt);
  \draw[newrung,fill=white] (bn) circle (2.2pt);
  \node[above=4pt] at (am) {$x_{i-1}$};
  \node[above=4pt] at (ai) {$x_i$};
  \node[above=4pt] at (ap) {$x_{i+1}$};
  \node[below=5pt] at (bm) {$b_{i-1}$};
  \node[below=5pt] at (bn) {$b_i'$};
  \node[below=5pt] at (bo) {$b_i$};
  \node[below=5pt] at (bp) {$b_{i+1}$};
  \node[left=18pt] at (am) {$\alpha$};
  \node[left=18pt] at (bm) {$\beta$};
  \draw[decorate,decoration={brace,amplitude=4pt}]
    (0,2.76)--node[above=5pt] {length in $[L,2L]$} (5.4,2.76);
  \draw[decorate,decoration={brace,amplitude=4pt}]
    (5.4,2.76)--node[above=5pt] {length in $[L,2L]$} (10.8,2.76);
\end{tikzpicture}
\endgroup
\caption{Narrowing a ladder in a single chart. Both displayed guide
  pieces read as geodesics. The interior mark $x_i$ is far from the
  end rungs, so quadrangle slimness supplies a new mate $b_i'$ on
  the opposite side, within $2\Delta$ of $x_i$. The old scheduled
  mate $b_i$ is at distance at most $h+2\Delta$ from $b_i'$ along
  that side. The drawing is
  schematic: the long side intervals and the two rung bounds use
  different visual scales.}
\label{fig:path-ladder}
\end{figure}
